%%%%%%%%%%%%%%%%%%%%%%%%%%%%%%%%%%%%%%%%%
% Short Sectioned Assignment LaTeX Template Version 1.0 (5/5/12)
% This template has been downloaded from: http://www.LaTeXTemplates.com
% Original author:  Frits Wenneker (http://www.howtotex.com)
% License: CC BY-NC-SA 3.0 (http://creativecommons.org/licenses/by-nc-sa/3.0/)
%%%%%%%%%%%%%%%%%%%%%%%%%%%%%%%%%%%%%%%%%

% \documentclass[paper=a4, fontsize=11pt]{scrartcl} % A4 paper and 11pt font size
\documentclass[11pt, a4paper]{book}
\usepackage[T1]{fontenc} % Use 8-bit encoding that has 256 glyphs
\usepackage[utf8]{inputenc}
\usepackage{fourier} % Use the Adobe Utopia font for the document - comment this line to return to the LaTeX default
\usepackage{listings} % para insertar código con formato similar al editor
\usepackage[spanish, es-tabla]{babel} % Selecciona el español para palabras introducidas automáticamente, p.ej. "septiembre" en la fecha y especifica que se use la palabra Tabla en vez de Cuadro
\usepackage{url} % ,href} %para incluir URLs e hipervínculos dentro del texto (aunque hay que instalar href)
\usepackage{graphics,graphicx, float} %para incluir imágenes y colocarlas
\usepackage[gen]{eurosym} %para incluir el símbolo del euro
\usepackage{cite} %para incluir citas del archivo <nombre>.bib
\usepackage{enumerate}
\usepackage{hyperref}
\usepackage{graphicx}
\usepackage{tabularx}
\usepackage{booktabs}
\usepackage{longtable}
\usepackage{lscape}

\usepackage[table,xcdraw]{xcolor}
\hypersetup{
	colorlinks=true,	% false: boxed links; true: colored links
	linkcolor=black,	% color of internal links
	urlcolor=cyan		% color of external links
}
\renewcommand{\familydefault}{\sfdefault}
\usepackage{fancyhdr} % Custom headers and footers
\pagestyle{fancyplain} % Makes all pages in the document conform to the custom headers and footers
\fancyhead[L]{} % Empty left header
\fancyhead[C]{} % Empty center header
\fancyhead[R]{Daniel González Serrano} % My name
\fancyfoot[L]{} % Empty left footer
\fancyfoot[C]{} % Empty center footer
\fancyfoot[R]{\thepage} % Page numbering for right footer
%\renewcommand{\headrulewidth}{0pt} % Remove header underlines
\renewcommand{\footrulewidth}{0pt} % Remove footer underlines
\setlength{\headheight}{13.6pt} % Customize the height of the header

\usepackage{titlesec, blindtext, color}
\definecolor{gray75}{gray}{0.75}
\newcommand{\hsp}{\hspace{20pt}}
\titleformat{\chapter}[hang]{\Huge\bfseries}{\thechapter\hsp\textcolor{gray75}{|}\hsp}{0pt}{\Huge\bfseries}
\setcounter{secnumdepth}{4}
\usepackage[Lenny]{fncychap}


\begin{document}

	% Plantilla portada UGR
	\begin{titlepage}
\newlength{\centeroffset}
\setlength{\centeroffset}{-0.5\oddsidemargin}
\addtolength{\centeroffset}{0.5\evensidemargin}
\thispagestyle{empty}

\noindent\hspace*{\centeroffset}\begin{minipage}{\textwidth}

\centering
\includegraphics[width=0.9\textwidth]{logos/logo_ugr.jpg}\\[1.4cm]

\textsc{ \Large TRABAJO FIN DE GRADO\\[0.2 cm]}
\textsc{ GRADO EN INGENIERÍA INFORMÁTICA}\\[1 cm]

{\Huge\bfseries Memes para todos \\}
\noindent\rule[-1ex]{\textwidth}{3 pt}\\[3.5 ex]
{\large\bfseries Almacenamiento Online }
\end{minipage}

\vspace{2.5cm}
\noindent\hspace*{\centeroffset}
\begin{minipage}{\textwidth}
\centering

\textbf{Autor}\\ {Daniel González Serrano}\\[2.5 ex]
\textbf{Director}\\ {Juan Julián Merelo Guervós}\\[2 cm]
\includegraphics[width=0.3\textwidth]{logos/etsiit_logo.png}\\[0.1 cm]
\textsc{Escuela Técnica Superior de Ingenierías Informática y de Telecomunicación}\\
\textsc{---}\\
Granada, \today
% Mostrar nombre de mes y año de entrega del documento

\end{minipage}
\end{titlepage}


	% Plantilla prefacio UGR
	\thispagestyle{empty}

\begin{center}
{\large\bfseries Memes para todos \\ Almacenamiento Online }\\
\end{center}
\begin{center}
Daniel González Serrano\\
\end{center}

%\vspace{0.7cm}

\vspace{0.5cm}
\noindent\textbf{Palabras clave}: \textit{software libre}
\vspace{0.7cm}

\noindent\textbf{Resumen}\\
	

\cleardoublepage

\begin{center}
	{\large\bfseries Igual, pero en inglés}\\
\end{center}
\begin{center}
	Daniel González Serrano\\
\end{center}
\vspace{0.5cm}
\noindent\textbf{Keywords}: \textit{open source}, \textit{floss}
\vspace{0.7cm}

\noindent\textbf{Abstract}\\


\cleardoublepage

\thispagestyle{empty}

\noindent\rule[-1ex]{\textwidth}{2pt}\\[4.5ex]

D. \textbf{Juan Julián Merelo Guervós}, Profesor(a) del ...

\vspace{0.5cm}

\textbf{Informo:}

\vspace{0.5cm}

Que el presente trabajo, titulado \textit{\textbf{Chief}},
ha sido realizado bajo mi supervisión por \textbf{Estudiante}, y autorizo la defensa de dicho trabajo ante el tribunal
que corresponda.

\vspace{0.5cm}

Y para que conste, expiden y firman el presente informe en Granada a \today

\vspace{1cm}

\textbf{El/la director(a)/es: }

\vspace{5cm}

\noindent \textbf{(nombre completo tutor/a/es)}

\chapter*{Agradecimientos}






	% Índice de contenidos
	\newpage
	\tableofcontents

	% Índice de imágenes y tablas
	\newpage
	\listoffigures

	% Si hay suficientes se incluirá dicho índice
	\listoftables 
	\newpage

  % Introducción 
	\chapter{Introducción}

¿Alguna vez ha tenido problemas a la hora de encontrar una imagen en su galería?

Si la respuesta es sí y que después de estar guardando contenido en su solución de almacenamiento elegida durante un tiempo, se vuelve insostenible y ya no es nada útil, el problema que se propone resolver este TFG le va a ser de interés. Al igual que usted, hay mucha gente que, debido a la gran cantidad de contenido digital que maneja le es imposible realizar una acción tan fácil (o difícil) como es la de encontrar una imagen.

\section{Contexto y motivación}

En este escenario de generación de contenido digital sin parangón se presenta un desafío para los creadores de contenido: la gestión y almacenamiento eficiente, eficaz
y ágil de dicho contenido. A pesar de contar con diversas plataformas de almacenamiento en la nube y soluciones de gestión de contenidos digitales, la mayoría de
estas soluciones tienen diferentes limitaciones que afectan a la organización, ordenación y accesibilidad de este contenido. La necesidad de una solución que permita
a los usuarios almacenar, organizar y acceder a su amplio repertorio contenido digital de forma eficiente, eficaz y ágil es evidente.

Esta falta de soluciones eficientes, eficaces y ágiles para la gestión y almacenamiento de contenido digital es la motivación principal de este proyecto. Aquí se
presenta una oportunidad significativa para desarrollar un sistema más robusto y adaptado a las necesidades actuales.

Al mejorar la eficiencia, eficacia y agilidad a la hora de la gestión y almacenamiento del contenido digital los usuarios pueden dedicarse a lo que realmente les importa:
la creación de contenido permitiéndoles así dedicarle más tiempo a la creatividad y a la expresión que requiere la creación de contenido digital sin miedo a perder o
no poder acceder a su contenido.

En este contexto y motivación, el presente proyecto se propone abordar estas limitaciones, diseñando y desarrollando un sistema de almacenamiento online que no solo almacene, 
sino que también organice, ordene y facilite el acceso a contenido digital de manera intuitiva y eficiente.

\section{Definición del problema}

Como se ha señalado previamente, la generación masiva de contenido digital ha creado un problema de gestión y almacenamiento del mismo. A pesar de la existencia de
soluciones en la nube y plataformas de gestión de contenido, existen limitaciones a la hora de la gestión y organización de imágenes. La ausencia de una solución que
permita la gestión y organización de imágenes de forma centralizada y eficiente es el problema que se pretende abordar con este proyecto.

    \subsection{Usuarios identificados}

    \begin{itemize}
        \item \textbf{Nombre}: María Rodríguez.
        \item \textbf{Edad}: 30 años.
        \item \textbf{Educación}: Licenciada en comunicación digital.
        \item \textbf{Ocupación}: Creadora de contenido.
        \item \textbf{Conocimientos y actitud hacia la tecnología}: Tiene conocimientos avanzados en herramientas digitales y plataformas de creación de contenido. Es entusiasta y proactiva hacia las nuevas tecnologías, siempre dispuesta a explorar nuevas soluciones que mejoren su experiencia y flujo de trabajo a la hora de la creación y gestión de contenido digital.
        \item \textbf{Dispositivos}:
            \begin{itemize}

            \item \textbf{¿Cuáles usa?}: Ordenador portátil de última generación, teléfono móvil de última generación y tablet.
            \item \textbf{¿Cuándo los usa?}: Durante todo el día.
            \item \textbf{¿Para qué los usa?}: Para mantenerse al tanto de las tendencias, interactuar con su audiencia y crear contenido.

            \end{itemize}
        \item \textbf{Misión}:
        \item 
            \begin{itemize}

                \item \textbf{¿Utilizaría nuestra aplicación?}: Sí.

                \item \textbf{¿Qué le gustaría que tuviera?}: Un sistema de almacenamiento en la nube que le permita subir y organizar fácilmente su contenido desde cualquier lugar. Le gustaría tener un acceso rápido y sencillo desde cualquier dispositivo con conexión a Internet, una interfaz intuitiva, un sistema de gestión de versiones, seguridad y privacidad, un sistema de búsqueda avanzada que permita encontrar cualquier contenido.

            \end{itemize}
        
    \end{itemize}

    \begin{itemize}
        \item \textbf{Nombre}: Juan Pérez.
        \item \textbf{Edad}: 40 años.
        \item \textbf{Educación}: Ingeniero de Sistemas.
        \item \textbf{Ocupación}: CTO.
        \item \textbf{Conocimientos y actitud hacia la tecnología}: Conocimientos avanzados en sistemas y tecnologías. Actitud positiva y pragmática.
        \item \textbf{Dispositivos}:
            \begin{itemize}
            \item \textbf{¿Cuáles usa?}: Ordenador de sobremesa y teléfono móvil. 
            \item \textbf{¿Cuándo los usa?}: Principalmente durante las horas laborales.
            \item \textbf{¿Para qué los usa?}: Para trabajar, informarse y comunicarse.
            \end{itemize}
        \item \textbf{Misión}:
            \begin{itemize}

                \item \textbf{¿Utilizaría nuestra aplicación?}: No. Prefiere un almacenamiento local, no considera necesario un sistema complejo y desconfía de la nube. Utiliza Microsoft OneDrive que ya le brinda la flexibilidad de acceso a los archivos, seguridad, privacidad y el sistema de búsqueda y gestión.

                \item \textbf{¿Qué le gustaría que tuviera?}: Le gustaría que no tuviera complicaciones a la hora de la instalación y configuración. Todas las facilidades que ya le brinda Microsoft OneDrive más alguna característica única e interesante para él.

            \end{itemize}
        
    \end{itemize}

    \begin{itemize}
        \item \textbf{Nombre}: Laura Martínez.
        \item \textbf{Edad}: 35 años.
        \item \textbf{Educación}: Licenciada en Bellas Artes.
        \item \textbf{Ocupación}: Fotógrafa profesional.
        \item \textbf{Conocimientos y actitud hacia la tecnología}: Tiene conocimientos avanzados en técnicas fotográficas y edición de imágenes. Prefiere soluciones sencillas y fáciles de usar. No le gusta perder el tiempo en configuraciones y opciones complejas.
        \item \textbf{Dispositivos}:
            \begin{itemize}

            \item \textbf{¿Cuáles usa?}: Ordenador portátil de alta gama y dispositivo móvil inteligente.
            \item \textbf{¿Cuándo los usa?}: Durante todo el día.
            \item \textbf{¿Para qué los usa?}: Revisar correos, redes sociales además de editar y visualizar imágenes.

            \end{itemize}
        \item \textbf{Misión}:
        \item 
            \begin{itemize}

                \item \textbf{¿Utilizaría nuestra aplicación?}: No. Siempre ha experimentado una pérdida de calidad al subir su contenido a la nube. Prefiere tener un almacenamiento local como el NAS que tiene para almacenar y organizar su contenido.

                \item \textbf{¿Qué le gustaría que tuviera?}: Su sistema idea es un sistema que facilite la organización, gestión y almacenamiento de su contenido de forma intuitiva, simple y eficaz sin pérdidas de calidad.

            \end{itemize}
        
    \end{itemize}

    \begin{itemize}
        \item \textbf{Nombre}: Carlos Sánchez
        \item \textbf{Edad}: 45 años.
        \item \textbf{Educación}: Licenciado en administración de empresas.
        \item \textbf{Ocupación}: Dueño de un pequeño negocio.
        \item \textbf{Conocimientos y actitud hacia la tecnología}: Conocimientos muy básicos en tecnología. Utiliza la tecnología para ganar difusión en su negocio aunque no sea un experto. Actitud conservadora y pragmática.
        \item \textbf{Dispositivos}:
            \begin{itemize}

            \item \textbf{¿Cuáles usa?}: Ordenador de sobremesa y teléfono inteligente.
            \item \textbf{¿Cuándo los usa?}: Durante las horas laborales.
            \item \textbf{¿Para qué los usa?}: Para atender llamadas y mensajes de clientes, comprobar el correo, redes sociales y gestionar su negocio.

            \end{itemize}
        \item \textbf{Misión}:
        \item 
            \begin{itemize}

                \item \textbf{¿Utilizaría nuestra aplicación?}: Sí, pero al tener una conexión inestable (a veces no funciona Internet) duda de si es una buena solución para él.

                \item \textbf{¿Qué le gustaría que tuviera?}: Valoraría mucho que el sistema ofreciera opciones de gestión offline y que no dependa de tener una conexión estable a Internet además de que fuera fácil de implementar y usar.

            \end{itemize}
        
    \end{itemize}

    \begin{itemize}
        \item \textbf{Nombre}: Marta González.
        \item \textbf{Edad}: 28 años.
        \item \textbf{Educación}: Graduada en Diseño Multimedia.
        \item \textbf{Ocupación}: Diseñadora Autónoma.
        \item \textbf{Conocimientos y actitud hacia la tecnología}: Tiene conocimientos avanzados en software de diseño gráfico y multimedia. Tiene una actitud positiva hacia la tecnología para mejorar su eficiencia en la organización y archivos digitales.
        \item \textbf{Dispositivos}:
            \begin{itemize}

            \item \textbf{¿Cuáles usa?}: Ordenador de sobremesa potente, teléfono inteligente y tableta gráfica.
            \item \textbf{¿Cuándo los usa?}: Durante las horas diurnas pero también algunas veces por la noche.
            \item \textbf{¿Para qué los usa?}: Para realizar sus proyectos de ilustración y diseño.

            \end{itemize}
        \item \textbf{Misión}:
        \item 
            \begin{itemize}

                \item \textbf{¿Utilizaría nuestra aplicación?}: Sí. Aunque siempre que ha utilizado un sistema para la organización de contenido ha tenido problemas cuando la cantidad del mismo es muy grande.

                \item \textbf{¿Qué le gustaría que tuviera?}: Le gustaría que fuera realmente escalable y mantenible aun con cantidades grandes de contenido.

            \end{itemize}
        
    \end{itemize}


\section{Historias de usuario}

\subsection{María Rodríguez (Creadora de contenido)}

\begin{enumerate}
    \item Como María Rodríguez, quiero poder acceder a la aplicación desde mi ordenador portátil, teléfono móvil y tablet en cualquier momento.
    \item Como María Rodríguez, quiero que la interfaz sea intuitiva que me permita navegar y gestionar mi galería de memes de forma rápida y sencilla, así no tengo que dedicar mucho tiempo a aprender a manejarla.
    \item Como María Rodríguez, quiero que mi galería de memes esté protegida para que no me roben contenido.
    \item Como María Rodríguez, me gustaría que existiera la posibilidad de volver recuperar la versión anterior de algún meme.
    \item Como María Rodríguez, necesito algún tipo de búsqueda avanzada para encontrar cualquier meme ya sea por título, etiquetas o contenido del propio meme.
\end{enumerate}

\subsection{Juan Pérez (CTO)}

    \begin{enumerate}
        \item Como Juan Pérez, quiero que la instalación y configuración se puedan hacer de manera sencilla y sin complicaciones.
        \item Como Juan Pérez, quiero que existan medidas de protección como lo hace Microsoft OneDrive.
        \item Como Juan Pérez, quiero que exista la posibilidad de disponer de una herramienta para crear memes de manera rápida y sencilla.
        \item Como Juan Pérez, quiero que la gestión de los memes se haga del modo más eficiente y rápida posible así los puedo encontrar más rápido.
    \end{enumerate}

\subsection{Laura Martínez (Fotógrafa Profesional)}

    \begin{enumerate}
        \item Como Laura Martínez, quiero un sistema de almacenamiento local como el NAS que ya tengo instalado que me permita mantener la calidad de mis imágenes y evitar la pérdida de calidad.
    \end{enumerate}

\subsection{Carlos Sánchez (Dueño de un pequeño negocio)}

    \begin{enumerate}
        \item Como Carlos Sánchez, prefiero una solución que ofrezca opciones de gestión aun sin conexión a internet.
        \item Como Carlos Sánchez, necesito una solución que ofrezca opciones de edición de memes aun sin conexión a internet.
        \item Como Carlos Sánchez, quiero que se permita compartir los memes fácilmente con las diferentes redes sociales para que tenga la opción de subir en cualquier momento cualquier meme relacionado con mi negocio para ganar popularidad.
        \item Como Carlos Sánchez, quiero poder acceder desde mi móvil y ordenador personal para, por ejemplo, crear y editar los memes en el ordenador personal y luego en el móvil subirlo a RRSS.
    \end{enumerate}

\subsection{Marta González (Diseñadora Freelancer)}

    \begin{enumerate}
        \item Como Marta González, quiero que sea fácil de mantener a medida que mi colección de memes crece.
        \item Como Marta González, me gustaría que la aplicación ofreciera integraciones con otras herramientas de diseño gráfico para facilitar la importación y exportación de contenido y poder conectar varias plataformas, aplicaciones, etc.
        \item Como Marta González, valoro mucho que se puedan manejar grandes cantidades de memes sin experimentar problemas de rendimiento o dificultades en la organización u ordenación.
    \end{enumerate}

\section{Objetivos}

La esencia de los objetivos radica en diseñar y desarrollar un sistema de almacenamiento en la nube. Este sistema se enfoca en resolver los diferentes problemas
identificados durante el análisis del problema. 

La experiencia de usuario es un pilar fundamental de nuestro proyecto. Varios de los problemas y de las personas identificadas en el análisis del problema
pueden resolverse con una útil, accesible, intuitiva y simple interfaz de usuario reduciendo los clics necesarios y disminuyendo el tiempo de carga,
todo con el objetivo de aumentar la satisfacción del usuario.

La gestión, ordenación y organización del contenido digital del usuario debe ser nuestro principal objetivo. Todo esto para reducir el tiempo que el usuario
dedica a la gestión de su contenido y aumentar el tiempo que dedica a la creación de contenido.

La compatibilidad con diferentes dispositivos y sistemas operativos es otro de los objetivos de este proyecto. El usuario debe poder acceder a su contenido
desde cualquier dispositivo con conexión a Internet y con cualquier sistema operativo.

La seguridad y privacidad del contenido del usuario es otro de los objetivos de este proyecto. El usuario debe poder confiar en el sistema de almacenamiento
y no debe preocuparse por la seguridad y privacidad de su contenido.

Además de estos objetivos existen otros como la incorporación de gestión de versiones para mejorar la experiencia del usuario y la implementación de un sistema
de búsqueda avanzada que permita al usuario encontrar cualquier contenido de forma rápida y sencilla.

Finalmente, otro objetivo de este proyecto es relativo a la evaluación y validación del sistema por parte del cliente, que, en nuestro caso es el tutor
y el tribunal. El sistema debe ser evaluado y validado por el cliente para comprobar que cumple con los requisitos y objetivos establecidos.

	\newpage

	\bibliography{bibliografia}
	\bibliographystyle{plain}
	
\end{document}

